\bta{2010}

\begin{enumerate}
	%\renewcommand{\labelenumi}{\arabic{enumi}.}
	% A(\Alph) a(\alph) I(\Roman) i(\roman) 1(\arabic)
	%设定全局标号series=example	%引用全局变量resume=example
	%[topsep=-0.3em,parsep=-0.3em,itemsep=-0.3em,partopsep=-0.3em]
	%可使用leftmargin调整列表环境左边的空白长度 [leftmargin=0em]
	\item
	选择题 (本题满分 40 分，每小题 5 分)
	
\begin{enumerate}
%\renewcommand{\labelenumi}{\arabic{enumi}.}
% A(\Alph) a(\alph) I(\Roman) i(\roman) 1(\arabic)
%设定全局标号series=example	%引用全局变量resume=example
%[topsep=-0.3em,parsep=-0.3em,itemsep=-0.3em,partopsep=-0.3em]
%可使用leftmargin调整列表环境左边的空白长度 [leftmargin=0em]
\item
已知 $\lim\limits _{x \rightarrow 0} \frac{a \tan x+b(1-\cos x)}{c \ln (1-2 x)+d\left(1-e^{-x^{2}}\right)}=2$, 其中 $a^{2}+c^{2}\neq 0$, 则必有 \xzanswer{D}

	
\fourchoices
{$b=4 d$}
{$b=-4 d$}
{$a=4 c$}
{$a=-4 c$}

\item 
$y=y(x)$ 满足 $\Delta y=\frac{1-x}{\sqrt{2 x-x^{2}}} \Delta x+o(\Delta x), y(1)=1, \int_{0}^{1} y(x)\,dx=$ \xzanswer{D} 


\fourchoices
{$2 \pi$}
{$\pi$}
{$\frac{\pi}{2}$}
{$\frac{\pi}{4}$}

\item 
设函数 $f(x)=\int_{0}^{x} t^{2}(t-1) d t$, 则 $f(x)$ 有多少个极值点 \xzanswer{B} 


\fourchoices
{$ 0 $}
{$ 1 $}
{$ 2 $}
{$ 3 $}


\item 
下列命题正确的一项是 \xzanswer{C} 


\fourchoices
{若 $\sum\limits_{n=1}^{\infty} a_{n}, \sum\limits_{n=1}^{\infty} b_{n}$ 都收敛，则 $\sum\limits_{n=1}^{\infty} a_{n} b_{n}$ 收敛}
{若 $\sum\limits_{n=1}^{\infty} a_{n}$ 收敛, $\sum\limits_{n=1}^{\infty} b_{n}$ 发散，则 $\sum\limits_{n=1}^{\infty} a_{n} b_{n}$ 发散}
{若 $\sum\limits_{n=1}^{\infty} a_{n}$ 收敛, $\sum\limits_{n=1}^{\infty} b_{n}$ 绝对收敛，则 $\sum\limits_{n=1}^{\infty} a_{n} b_{n}$ 绝对收敛}
{若 $\sum\limits_{n=1}^{\infty} a_{n}$ 收敛, $\sum\limits_{n=1}^{\infty} b_{n}$ 条件收敛，则 $\sum\limits_{n=1}^{\infty} a_{n} b_{n}$ 条件收敛}


\item
设 $a_{0}, a_{1}, a_{2},a_{3} \ldots$ 是等差数列, 且公差 $d>0$, 则幂级数 $\sum\limits_{n=0}^{\infty} a_{n} x^{n}$ 的收敛域为 \xzanswer{C}

	

\fourchoices
{$(-d, d)$}
{$[-d, d)$}
{$(-1,1)$}
{$[-1,1)$}


\item 
设 $y_{1}(x), y_{2}(x), y_{3}(x)$ 是二阶线性非齐次微分方程 $y^{\prime \prime}+p(x) y^{\prime}+q(x) y=f(x)$ 的三 个线性无关的解， $C_{1}, C_{2}$ 是两个任意常数, 则微分方程的通解为 \xzanswer{B} 

\fourchoices
{$C_{1} y_{1}(x)+C_{2} y_{2}(x)+C_{3} y_{3}(x)$}
{$C_{1} y_{1}(x)+C_{2} y_{2}(x)+\left(1-C_{1}-C_{2}\right) y_{3}(x)$}
{$C_{1} y_{1}(x)+C_{2} y_{2}(x)-\left(C_{1}+C_{2}\right) y_{3}(x)$}
{$C_{1} y_{1}(x)+C_{2} y_{2}(x)-\left(1-C_{1}-C_{2}\right) y_{3}(x)$}

\item 
通过两个平面 $2 x+y-4=0$ 与 $y+2 z=0$ 的交线及点 $M_{0}(2,-1,-1)$ 的平面方程为 \xzanswer{A} 


\fourchoices
{$3 x+y-z=6$}
{$x+3 y-z=0$}
{$3 x-y+z=6$}
{$x-3 y-z=6$}

\item 
曲线 $y=e^{x}$ 和该曲线经过原点的切线以及 $y$ 轴所围成的面积 \xzanswer{A} 


\fourchoices
{$\frac{e}{2}-1$}
{$\frac{e}{2}+1$}
{$\frac{e}{2}$}
{$e+1$}



\end{enumerate}


\banswer{
\begin{enumerate}
	%\renewcommand{\labelenumi}{\arabic{enumi}.}
	% A(\Alph) a(\alph) I(\Roman) i(\roman) 1(\arabic)
	\item
	\textbf{D}。展开到二阶：分子
	$a\tan x+b(1-\cos x)=ax+\frac b2x^2+o\left(x^2\right)$，分母
	$c\ln(1-2x)+d\left(1-e^{-x^2}\right)=-2cx+(d-2c)x^2+o\left(x^2\right)$。

	若 $c=0$（则 $a\neq0$），分母是 $x^2$ 阶而分子是 $x$ 阶，极限为 $\infty$，不合；
	若 $a=0$（则 $c\neq0$），分子是 $x^2$ 阶而分母是 $x$ 阶，极限为 $0\neq2$，不合。
	故 $a\neq0$ 且 $c\neq0$，此时
	\[
	\lim\limits_{x\to0}\frac{ax+o(x)}{-2cx+o(x)}=-\frac{a}{2c}=2
	\quad\Longrightarrow\quad a=-4c .
	\]

	\item
	\textbf{D}。$\Delta y=y'\Delta x+o(\Delta x)$ 说明
	$y'=\frac{1-x}{\sqrt{2x-x^2}}$。注意 $\left(2x-x^2\right)'=2(1-x)$，故
	\[
	y=\int\frac{\frac12 d\left(2x-x^2\right)}{\sqrt{2x-x^2}}=\sqrt{2x-x^2}+C ,
	\]
	由 $y(1)=1$ 得 $C=0$，即 $y=\sqrt{1-(x-1)^2}$。
	它在 $[0,1]$ 上的图形是单位圆的四分之一弧，故
	$\int_0^1y\,dx=\frac{\pi}{4}$。

	\item
	\textbf{B}。$f'(x)=x^{2}(x-1)$。$x=0$ 是二重根，两侧 $f'$ 同号（均为负），不是极值点；
	$x=1$ 处 $f'$ 由负变正，是极小值点。故只有 $1$ 个极值点。

	\item
	\textbf{C}。$\sum a_n$ 收敛 $\Rightarrow a_n\to0$，从而 $|a_n|\leqslant M$；
	于是 $\sum|a_nb_n|\leqslant M\sum|b_n|<\infty$，即 $\sum a_nb_n$ 绝对收敛。

	反例排除其余：
	A 错，取 $a_n=b_n=\frac{(-1)^n}{\sqrt n}$，两级数收敛而 $\sum a_nb_n=\sum\frac1n$ 发散；
	B 错，取 $a_n=\frac1{n^2}$、$b_n=1$，则 $\sum a_nb_n$ 收敛；
	D 错，取 $a_n=\frac1{n^2}$、$b_n=\frac{(-1)^n}{n}$，则 $\sum a_nb_n=\sum\frac{(-1)^n}{n^3}$
	绝对收敛，并非条件收敛。

	\item
	\textbf{C}。$a_n=a_0+nd$ 随 $n$ 线性增长，$\sqrt[n]{|a_n|}\to1$，
	由柯西--阿达马公式收敛半径 $R=1$。
	端点 $x=\pm1$ 处通项 $\left|a_n(\pm1)^n\right|=|a_n|\to+\infty$ 不趋于 $0$，
	级数发散。故收敛域为 $(-1,1)$。

	\item
	\textbf{B}。$y_1-y_3$ 与 $y_2-y_3$ 是齐次方程的两个线性无关解，
	故通解为
	\[
	y_3+C_1\left(y_1-y_3\right)+C_2\left(y_2-y_3\right),
	\]
	即 $C_1y_1+C_2y_2+\left(1-C_1-C_2\right)y_3$。
	（判别法：非齐次方程解的线性组合仍为解的充要条件是系数之和为 $1$；
	A 多一个常数，C、D 的系数和分别为 $0$ 和 $2C_1+2C_2-1$，皆不恒为 $1$。）

	\item
	\textbf{A}。设过交线的平面束 $2x+y-4+\lambda(y+2z)=0$。
	代入 $M_0(2,-1,-1)$ 得 $-1-3\lambda=0$，$\lambda=-\frac13$，整理得
	\[
	6x+2y-2z-12=0,\qquad\text{即}\qquad 3x+y-z=6 .
	\]

	\item
	\textbf{A}。设切点为 $\left(a,e^{a}\right)$，切线 $y=e^{a}(x-a)+e^{a}$ 过原点给出
	$e^{a}(1-a)=0$，即 $a=1$，切线为 $y=ex$。所求面积
	\[
	S=\int_0^1\left(e^{x}-ex\right)dx=(e-1)-\frac e2=\frac e2-1 .
	\]
\end{enumerate}
}


\item
(本题满分 12 分)
设 $g(x)$ 是以正数 $T$ 为周期的连续函数, $g(0)=1, f(x)=\int_{0}^{2 x}|x-t| g(t) d t$, 求 $f^{\prime}(T)$ 。

\banswer{
	\textbf{去绝对值.} 设 $x>0$，按 $t=x$ 分段：
	\begin{align*}
	f(x)&=\int_0^x\!(x-t)g(t)dt+\int_x^{2x}\!(t-x)g(t)dt\\
	&=x\int_0^x\!g-\int_0^x\!tg\,dt+\int_x^{2x}\!tg\,dt-x\int_x^{2x}\!g .
	\end{align*}

	\textbf{求导.} 逐项使用变限积分求导（注意上限 $2x$ 的链式法则）：
	\begin{align*}
	\left(x\int_0^x\!g\right)'&=\int_0^x\!g+xg(x),\\
	\left(-\int_0^x\!tg\right)'&=-xg(x),\\
	\left(\int_x^{2x}\!tg\right)'&=4xg(2x)-xg(x),\\
	\left(-x\int_x^{2x}\!g\right)'&=-\int_x^{2x}\!g-x\left[2g(2x)-g(x)\right].
	\end{align*}
	相加，含 $g(x)$、$g(2x)$ 的项大量抵消，得
	\[
	f'(x)=\int_0^x\!g(t)dt-\int_x^{2x}\!g(t)dt+2xg(2x) .
	\]

	\textbf{取 $x=T$.} 由周期性 $\int_T^{2T}g=\int_0^Tg$，前两项抵消；
	又 $g(2T)=g(0)=1$，故
	\[
	f'(T)=2T\cdot g(2T)=2T .
	\]
	（数值检验：取 $T=2$ 及若干个周期为 $2$、$g(0)=1$ 的函数，
	差商均给出 $f'(T)=4.0000$。）
}



\item 
(本题满分 12 分)
已知平面过点 $ (1,2,3) $，它在 $ x $ 正半轴和 $ y $ 轴正半轴上的截距相等，问当该平面在
三个坐标轴上的截距分别为何值时，它与三个坐标轴所围成的立体体积取极值，并写出此时的平面方程。


%群友反馈：原题目有问题，最小值为零，两种改法：①将$ z $轴的截距限制为正；②将最值改为极值。


\banswer{
	\textbf{建模.} 设 $x$、$y$ 轴上的截距同为 $a>0$，$z$ 轴上的截距为 $c$，
	则平面为 $\frac xa+\frac ya+\frac zc=1$。过 $(1,2,3)$ 给出约束
	\[
	\frac1a+\frac2a+\frac3c=1,\qquad\text{即}\qquad\frac3a+\frac3c=1 ,
	\]
	立体体积 $V=\frac16a\cdot a\cdot c=\frac{a^2c}{6}$。

	\textbf{消元求极值.} 由约束得 $c=\frac{3a}{a-3}$（取 $a>3$ 以保证 $c>0$），于是
	\[
	V(a)=\frac{a^{3}}{2(a-3)},\qquad
	V'(a)=\frac{a^{2}(2a-9)}{2(a-3)^{2}} .
	\]
	$V'$ 在 $a=\frac92$ 处由负变正，故取极小值；此时
	\[
	c=\frac{3\cdot\frac92}{\frac92-3}=9,\qquad V_{\min}=\frac{243}{8} .
	\]

	\textbf{结论.} $x$ 轴、$y$ 轴、$z$ 轴上的截距分别为 $\frac92$、$\frac92$、$9$ 时体积取极（小）值
	$\frac{243}{8}$，此时平面方程为
	\[
	\frac{x}{9/2}+\frac{y}{9/2}+\frac{z}{9}=1,\qquad\text{即}\qquad 2x+2y+z=9 .
	\]

	\par\vspace{1mm}\noindent
	\textbf{注：}{\kaishu
	原答案把 $x$、$y$ 截距写作 $\frac29$，与它自己给出的平面方程
	$\frac{x}{9/2}+\frac{y}{9/2}+\frac z9=1$ 自相矛盾，应为 $\frac92$；
	原答案又把“平面方程”误写成“曲面方程”。数值搜索确认 $a=\frac92$，$V=\frac{243}{8}=30.375$。

	题面写“取极值”而非“取最值”是必要的：若允许 $0<a<3$，则 $c<0$，
	体积 $|V|=\frac{a^3}{2(3-a)}$ 当 $a\to0^{+}$ 时趋于 $0$ 却取不到，全局最小值不存在。
	本题与 2011 年第十一题完全相同。
	}\par
}


\item 
(本题满分 12 分)
求下面微分方程的通解：
\[ 
y y^{\prime \prime}-\left(y^{\prime}\right)^{2}=y^{2} \ln y
 \]
注: $\int \frac{d x}{\sqrt{x^{2} + a^{2}}}=\ln \left(x+\sqrt{x^{2} + a^{2}}\right)$+C。


\banswer{
    原方程含 $\ln y$，故实解必须满足 $y>0$。令 $w=\ln y$，则
    \[
    w''=\frac{yy''-(y')^2}{y^2}=\ln y=w.
    \]
    所以 $w''-w=0$，其完整通解为 $w=C_1e^x+C_2e^{-x}$，从而
    \[
    \boxed{y(x)=\exp\!\left(C_1e^x+C_2e^{-x}\right)},\qquad C_1,C_2\in\mathbb R.
    \]
    回代可得 $yy''-(y')^2=y^2w''=y^2\ln y$。
    该表达式包含常数解 $y\equiv1$，也包含 $y'$ 正、负以及改变符号的分支；
    只取 $y'=y\sqrt{\ln^2y+C}$ 会漏掉解，不能称为完整通解。
}


\item 
(本题满分 12 分)
将函数 $f(x)=\frac{x-1}{(x+1)^{2}}$ 在 $x=0$ 处展开成幂级数, 并求收敛区间。


\banswer{
	\textbf{出发点.} 由 $\frac{1}{1+x}=\sum\limits_{n=0}^{\infty}(-1)^{n}x^{n}$（$|x|<1$）逐项求导并取负号，
	\[
	\frac{1}{(1+x)^{2}}=\sum\limits_{n=0}^{\infty}(-1)^{n}(n+1)x^{n},\qquad |x|<1 .
	\]

	\textbf{乘上 $x-1$.}
	\begin{align*}
	f(x)&=(x-1)\sum\limits_{n=0}^{\infty}(-1)^{n}(n+1)x^{n}\\
	&=\sum\limits_{n=1}^{\infty}(-1)^{n-1}nx^{n}-\sum\limits_{n=0}^{\infty}(-1)^{n}(n+1)x^{n} .
	\end{align*}
	合并 $x^{n}$ 的系数（$n\geqslant1$）：
	\[
	(-1)^{n-1}n-(-1)^{n}(n+1)=(-1)^{n+1}(2n+1) ,
	\]
	而 $n=0$ 项为 $-1$，恰与公式在 $n=0$ 的值一致。故
	\[
	f(x)=\sum\limits_{n=0}^{\infty}(-1)^{n+1}(2n+1)x^{n}
	=-1+3x-5x^{2}+7x^{3}-\cdots
	\]

	\textbf{收敛区间.} 系数满足 $\left|\frac{a_{n+1}}{a_n}\right|=\frac{2n+3}{2n+1}\to1$，
	故 $R=1$；端点 $x=\pm1$ 处通项不趋于 $0$，发散。收敛区间为 $(-1,1)$。
}


\item
(本题满分 12 分)
在椭球面 $2 x^{2}+2 y^{2}+z^{2}=1$ 上求一点, 使函数 $f(x, y, z)=x^{2}+y^{2}+z^{2}$ 在该点, 沿 方向 $\vec{l}=\vec{i}-\vec{j}$ 的方向导数最大。


\banswer{
	\textbf{写出方向导数.} $\nabla f=(2x,2y,2z)$，$\vec l=\vec i-\vec j$ 的单位向量为
	$\frac{1}{\sqrt2}(1,-1,0)$，故
	\[
	\frac{\partial f}{\partial\vec l}=\nabla f\cdot\frac{\vec l}{|\vec l|}
	=\frac{2x-2y}{\sqrt2}=\sqrt2\,(x-y) .
	\]
	问题化为：在 $2x^{2}+2y^{2}+z^{2}=1$ 上求 $x-y$ 的最大值。

	\textbf{求最大值.} 由柯西不等式
	\[
	(x-y)^{2}\leqslant\left(1^2+1^2\right)\left(x^{2}+y^{2}\right)
	=2\cdot\frac{2x^2+2y^2}{2}\leqslant 2\cdot\frac{1}{2}=1 ,
	\]
	故 $x-y\leqslant1$，且等号要求 $x=-y$、$z=0$ 且 $2x^2+2y^2=1$，
	即 $x=\frac12,\ y=-\frac12,\ z=0$。

	\textbf{结论.} 在点 $\left(\frac12,-\frac12,0\right)$ 处方向导数取最大值
	\[
	\left.\frac{\partial f}{\partial\vec l}\right|_{\max}=\sqrt2 .
	\]
	（也可用拉格朗日乘数法；在椭球面上取 $4\times10^{5}$ 个采样点数值搜索，
	最大值为 $1.41421$，取值点为 $(0.5,-0.5,0)$，与上吻合。）
}

\newpage
\item 
(本题满分 14 分)
设 $f(x)$ 为定义在 $[0,+\infty)$ 上的连续函数,且满足
\[ 
f(t)=\iiint\limits_{x^{2}+y^{2}+z^{2} \leqslant t^{2}} f\left(\sqrt{x^{2}+y^{2}+z^{2}}\right) d v+t^{3}
 \]
求 $f(x)$ 的表达式。

\banswer{
	\textbf{化三重积分为单积分.} 被积函数只依赖 $r=\sqrt{x^2+y^2+z^2}$，
	用球坐标（球壳体积元 $4\pi r^{2}dr$）得
	\[
	f(t)=4\pi\int_{0}^{t}r^{2}f(r)dr+t^{3},\qquad t\geqslant0 .
	\]

	\textbf{化为微分方程.} 令 $t=0$ 得 $f(0)=0$。由 $f$ 连续，右端可导，两端求导：
	\[
	f'(t)=4\pi t^{2}f(t)+3t^{2},\qquad\text{即}\qquad f'-4\pi t^{2}f=3t^{2} .
	\]

	\textbf{解方程.} 积分因子 $e^{-\frac{4\pi}{3}t^{3}}$，于是
	$\left(fe^{-\frac{4\pi}{3}t^{3}}\right)'=3t^{2}e^{-\frac{4\pi}{3}t^{3}}$，积分得
	\[
	f(t)e^{-\frac{4\pi}{3}t^{3}}=-\frac{3}{4\pi}e^{-\frac{4\pi}{3}t^{3}}+C .
	\]
	由 $f(0)=0$ 得 $C=\frac{3}{4\pi}$，故
	\[
	f(x)=\frac{3}{4\pi}\left(e^{\frac{4}{3}\pi x^{3}}-1\right),\qquad x\geqslant0 .
	\]
	（代回原积分方程，两端之差恒为 $0$，已用符号计算验证。）
}



\item 
(本题满分 12 分)
设 $u=u\left(\sqrt{x^{2}+y^{2}}\right)$ 具有连续的二阶偏导, 且满足 
\[ 
\frac{\partial^{2} u}{\partial x^{2}}+\frac{\partial^{2} u}{\partial y^{2}}-\frac{1}{x} \frac{\partial u}{\partial x}+u=x^{2}+y^{2}
 \]
 试求 $u$ 的表达式。

\banswer{
	\textbf{化为常微分方程.} 记 $r=\sqrt{x^{2}+y^{2}}$，$u=u(r)$。由
	$r_x=\frac xr$、$r_y=\frac yr$ 得
	\[
	\frac{\partial u}{\partial x}=u'(r)\frac xr,\qquad
	\frac{\partial^{2}u}{\partial x^{2}}+\frac{\partial^{2}u}{\partial y^{2}}
	=u''(r)+\frac{u'(r)}{r} .
	\]
	注意 $\frac1x\frac{\partial u}{\partial x}=\frac{u'(r)}{r}$，
	恰与拉普拉斯算子中的一阶项抵消，于是原方程化为
	\[
	u''(r)+u(r)=r^{2} .
	\]

	\textbf{解方程.} 齐次部分特征根 $\pm i$，给出 $c_1\cos r+c_2\sin r$；
	设特解 $u^{*}=Ar^{2}+Br+C$，代入得 $2A+Ar^{2}+Br+C=r^{2}$，
	比较系数 $A=1,\ B=0,\ C=-2$。故
	\[
	u(r)=c_{1}\cos r+c_{2}\sin r+r^{2}-2 ,
	\]
	即
	\[
	u=c_{1}\cos\sqrt{x^{2}+y^{2}}+c_{2}\sin\sqrt{x^{2}+y^{2}}+x^{2}+y^{2}-2 .
	\]

	\par\vspace{1mm}\noindent
	\textbf{注：}{\kaishu
	本题的关键是 $-\frac1x\frac{\partial u}{\partial x}$ 恰好把极坐标拉普拉斯算子里的
	$\frac1r\frac{du}{dr}$ 消去，方程才成为常系数的 $u''+u=r^{2}$。
	又因方程中出现 $\frac1x$，所求解只在 $x\neq0$ 的区域内有意义；
	若还要求 $u$ 作为 $(x,y)$ 的函数在原点附近光滑，则须取 $c_2=0$，
	因为 $\sin\sqrt{x^2+y^2}$ 在原点不可微。
	}\par
}

\item 
(本题满分 12 分)
设函数 $f(x)$ 具有连续导数，在区域 $G$ 内曲线积分
\[ 
\int_{M}^{N} \frac{1}{2 x^{2}+f(y)}(y d x-x d y)
 \]
与路径无关，其中 $G$ 不包含原点的单连通区域, $M 、 N$ 是 $G$ 内的任意两点, 且 $f(1)=1$ 。
\begin{enumerate}
	%\renewcommand{\labelenumi}{\arabic{enumi}.}
	% A(\Alph) a(\alph) I(\Roman) i(\roman) 1(\arabic)
	%设定全局标号series=example	%引用全局变量resume=example
	%[topsep=-0.3em,parsep=-0.3em,itemsep=-0.3em,partopsep=-0.3em]
	%可使用leftmargin调整列表环境左边的空白长度 [leftmargin=0em]
	\item
 求 $f(x)$;
\item 
求 $\oint_{\Gamma} \frac{1}{2 x^{2}+f(y)}(y d x-x d y)$ 其中 $\Gamma$ 为闭曲线 $x^{\frac{2}{3}}+y^{\frac{2}{3}}=a^{\frac{2}{3}}$ $ (a>0) $, 取逆时针方向。
\end{enumerate}

\banswer{
\begin{enumerate}
	%\renewcommand{\labelenumi}{\arabic{enumi}.}
	% A(\Alph) a(\alph) I(\Roman) i(\roman) 1(\arabic)
	%设定全局标号series=example	%引用全局变量resume=example
	%[topsep=-0.3em,parsep=-0.3em,itemsep=-0.3em,partopsep=-0.3em]
	%可使用leftmargin调整列表环境左边的空白长度 [leftmargin=0em]
	\item
	记 $P=\dfrac{y}{2x^{2}+f(y)}$，$Q=\dfrac{-x}{2x^{2}+f(y)}$。
	在单连通区域内曲线积分与路径无关的充要条件是 $\frac{\partial P}{\partial y}=\frac{\partial Q}{\partial x}$。
	直接求导（记 $S=2x^{2}+f(y)$）：
	\[
	\frac{\partial P}{\partial y}=\frac{S-yf'(y)}{S^{2}},\qquad
	\frac{\partial Q}{\partial x}=\frac{-S+4x^{2}}{S^{2}}=\frac{2x^{2}-f(y)}{S^{2}} .
	\]
	令两者相等，$2x^{2}+f(y)-yf'(y)=2x^{2}-f(y)$，即
	\[
	yf'(y)=2f(y)\quad\Longrightarrow\quad f(y)=Cy^{2} .
	\]
	由 $f(1)=1$ 得 $C=1$，故 $f(x)=x^{2}$。

	\item
	此时被积表达式为 $\dfrac{y\,dx-x\,dy}{2x^{2}+y^{2}}$，唯一奇点是原点，而星形线 $\Gamma$ 绕原点一周。
	取充分小的 $\varepsilon>0$，以椭圆
	$\ell_\varepsilon:\ 2x^{2}+y^{2}=\varepsilon^{2}$（逆时针）挖去原点；
	在 $\Gamma$ 与 $\ell_\varepsilon$ 之间的复连通区域上由 (1) 知 $P_y=Q_x$，
	故格林公式给出 $\oint_{\Gamma}=\oint_{\ell_\varepsilon}$。

	令 $x=\frac{\varepsilon}{\sqrt2}\cos t,\ y=\varepsilon\sin t$，则
	\[
	y\,dx-x\,dy=-\frac{\varepsilon^{2}}{\sqrt2}dt,\qquad 2x^{2}+y^{2}=\varepsilon^{2} ,
	\]
	于是
	\[
	\oint_{\Gamma}=\int_{0}^{2\pi}\left(-\frac{1}{\sqrt2}\right)dt=-\sqrt{2}\,\pi .
	\]
	（沿 $a=1$ 的星形线直接数值积分得 $-4.44288$，而 $-\sqrt2\pi=-4.44288$，吻合。）
\end{enumerate}


}


\item 
(本题满分 10 分)
设 $f(x)$ 在 $[0,1]$ 上具有二阶连续导数, 且 $f(0)=f(1)=0, f(x)$ 不恒为零, 求证:
$$
\int_{0}^{1}\left|f^{\prime \prime}(x)\right| d x \geqslant 4 \max _{0 \leqslant x \leqslant 1}|f(x)|
$$


\banswer{
令 $M=\max_{[0,1]}|f(x)|$。因 $f$ 连续、$f(0)=f(1)=0$ 且 $f\not\equiv0$，
存在 $x_0\in(0,1)$ 使 $|f(x_0)|=M>0$。不妨设 $f(x_0)=M$（否则考虑 $-f$）。
由拉格朗日中值定理，存在 $\xi\in(0,x_0)$，$\eta\in(x_0,1)$ 使得
\[
f'(\xi)=\frac{f(x_0)-f(0)}{x_0}=\frac{M}{x_0},\qquad
f'(\eta)=\frac{f(1)-f(x_0)}{1-x_0}=-\frac{M}{1-x_0}.
\]
于是
\[
\int_0^1|f''(x)|dx\ge\int_{\xi}^{\eta}|f''(x)|dx
\ge\Bigl|\int_{\xi}^{\eta}f''(x)dx\Bigr|
=|f'(\eta)-f'(\xi)|
=M\Bigl(\frac1{x_0}+\frac1{1-x_0}\Bigr).
\]
由 $x_0(1-x_0)\le\frac14$ 得 $\frac1{x_0}+\frac1{1-x_0}=\frac1{x_0(1-x_0)}\ge4$，
故 $\int_0^1|f''(x)|dx\ge4M$。
}


\item 
(本题满分 10 分)
一点从坐标原点出发向东移动 $ a \ \mathrm{m} $，然后左拐弯移动 $ aq \ \mathrm{m} $(其中 $ 0<q<1 $)，此后反
复左拐弯前行，使得后一段移动为前一段的 $ q $ 倍，该点这样运动下去，有一极限位置，
求该极限位置距离原点距离。


\banswer{
	\textbf{解法一（按坐标分组求和）.} 以原点为起点、正东为 $x$ 轴正向、正北为 $y$ 轴正向。
	每次左拐 $90^{\circ}$，故第 $k$ 段（$k=0,1,2,\dots$）长 $aq^{k}$，方向依次为
	东、北、西、南循环。把 $x$、$y$ 方向的位移分别求和：
	\begin{align*}
	x&=a-aq^{2}+aq^{4}-\cdots=\frac{a}{1+q^{2}},\\
	y&=aq-aq^{3}+aq^{5}-\cdots=\frac{aq}{1+q^{2}} ,
	\end{align*}
	两个几何级数的公比均为 $-q^{2}$，$|q|<1$ 故收敛。极限位置为
	$\left(\frac{a}{1+q^2},\frac{aq}{1+q^2}\right)$，到原点的距离
	\[
	d=\frac{a}{1+q^{2}}\sqrt{1+q^{2}}=\frac{a}{\sqrt{1+q^{2}}} .
	\]

	\textbf{解法二（每四步一组的向量迭代）.} 记 $\vv{n}_{0}=(1,q)$。
	前四步的合位移为
	\[
	a(1,0)+aq(0,1)+aq^{2}(-1,0)+aq^{3}(0,-1)=a\left(1-q^{2}\right)\vv{n}_{0} ,
	\]
	而此后每四步只是把上一组整体乘以 $q^{4}$。故总位移
	\[
	a\left(1-q^{2}\right)\vv{n}_{0}\sum_{k=0}^{\infty}q^{4k}
	=\frac{a\left(1-q^{2}\right)}{1-q^{4}}\vv{n}_{0}=\frac{a}{1+q^{2}}\vv{n}_{0} ,
	\]
	取模同样得到 $d=\frac{a}{\sqrt{1+q^{2}}}$。

	（数值检验：$q=0.3,0.5,0.8$ 时逐步累加 $400$ 段所得距离分别为
	$0.95783a$、$0.89443a$、$0.78087a$，与公式完全一致。）
}




\end{enumerate}

